\subsection{Introduction}

Le son est une onde mécanique qui se propage dans un gaz sous forme de variations de pression. Sa vitesse dépend des propriétés du milieu, notamment de sa température et de sa masse molaire, et elle ne dépend pas de la fréquence.

Dans ce laboratoire, le trombone de Koenig permet de mesurer la vitesse du son dans différents gaz. Cet appareil est un interféromètre pour les ondes sonores, un haut-parleur alimenté par un générateur de fréquence émet un son qui se propage selon deux chemins, l'un fixe et l'autre de longueur variable grâce à la partie mobile du tube. Un microphone relié à un oscilloscope capte la superposition des deux ondes. En déplaçant le tube mobile, on observe alternativement des interférences constructives et destructives, ce qui permet de déterminer la longueur d'onde $\lambda$ du son. La vitesse de propagation s'en déduit à l'aide de la fréquence $f$, mesurée avec précision par un fréquencemètre.

Nous avons réalisé les mesures décrites dans le descriptif PHY3-03-TPE12 \citeit{descriptif_tpe12} pour deux gaz, l'air et le dioxyde de carbone (CO$_2$), sur un domaine de fréquences que nous avons défini pour chaque milieu. Pour chaque gaz, la vitesse mesurée est obtenue par régression linéaire en tenant compte des incertitudes, puis comparée à la vitesse théorique calculée à partir du facteur adiabatique, de la température et de la masse molaire du gaz. Cette comparaison permet de vérifier que la vitesse du son est bien indépendante de la fréquence et qu'elle suit la loi attendue.